\documentclass[10pt]{article}
\usepackage[a4paper,margin=2cm]{geometry}
\usepackage[T1]{fontenc}
\usepackage{lmodern}
\usepackage{textcomp}
\usepackage{amsmath}
\usepackage{booktabs,tabularx,enumitem}
\usepackage[hidelinks]{hyperref}
\setlist{nosep,leftmargin=1.2em}
\setlength{\parskip}{3pt}
\setlength{\parindent}{0pt}
\renewcommand{\arraystretch}{1.15}
\newcolumntype{L}{>{\raggedright\arraybackslash}X}

\begin{document}

{\Large\bfseries Fairness over time in dynamic kidney exchange: a research summary}\\[2pt]
{\small Dr Yousef Amer (Adelaide University, Australia) and Neha (Manipal University Jaipur, India).
Prepared for potential collaborators.}

\section*{1. Overview}
Kidney exchange programmes let patients with a willing but incompatible living donor
swap donors with other pairs, through short cycles and through chains started by
altruistic donors. A programme clears its pool repeatedly: each month it selects a set
of exchanges, matched pairs leave, unmatched patients keep waiting, some withdraw, some
receive a deceased-donor kidney, and new pairs arrive. Fairness in kidney exchange,
however, is mostly studied within a single clearing round.

Our setting is the Australian and New Zealand Paired Kidney Exchange (ANZKX). Public
reports show 60--90 registered pairs at any time and 50--66 transplants per year, and
about 30\% of patients are highly sensitised (hard to match). We built a synthetic
simulator informed by this public information and a multi-objective framework that
trades transplant volume against three fairness criteria. Our central question is how
the price of fairness evolves across sequential clearing rounds under operational
constraints, and which clearing strategy, myopic or batched, better preserves the
balance between efficiency and equity over time. A related paper on the static setting
is under review, and a manuscript on the dynamic setting is in preparation.

\textbf{All results below come from synthetic simulation informed by public ANZKX
information. They are not clinical results and were not validated on patient-level data.}

\section*{2. Research gaps}
Six gaps in the literature motivate the work. Table~\ref{tab:gaps} lists each gap, how
we address it, and what we have found so far.

\begin{table}[h]
\small
\begin{tabularx}{\textwidth}{@{}p{0.5cm}p{3.1cm}LL@{}}
\toprule
 & Gap & How we address it & Finding so far (simulation) \\
\midrule
G1 & Static pool assumption: fairness-aware models treat the pool as fixed.
   & A dynamic pool over 25 monthly rounds with Poisson arrivals and withdrawal after
     $\tau \in \{6, 12, 18\}$ months.
   & Fairness gains persist over the whole trajectory, not only in one round. \\
G2 & Deceased-donor pathway left out of fairness-aware models.
   & Each waiting patient may leave for a deceased-donor transplant (competing exit,
     rate $\mu_d$).
   & Changing $\mu_d$ left the price of fairness unchanged but hurt batched clearing
     far more than myopic clearing ($-26\%$ vs.\ $-1.6\%$ transplants). \\
G3 & Logistics constraints absent: inter-hospital cold ischaemia time (CIT) is ignored.
   & Six regional nodes with door-to-door transport times and a CIT limit; logistics
     disparity $F_4$ as a modelled proxy.
   & A 12\,h CIT limit removed 13--14\% of transplants for every method. \\
G4 & Clearing strategies compared on volume only.
   & Myopic and batched clearing compared on all four objectives.
   & Monthly clearing transplanted more patients, including more highly sensitised
     patients, than every batching variant tested. \\
G5 & Price of fairness (PoF) measured only in single rounds.
   & The trajectory $\{\mathrm{PoF}_t\}_{t=1}^{T}$ and the whole-trajectory loss
     $\Delta$ against a volume-maximising policy.
   & PoF is about 2\% per round and does not accumulate; most deferred transplants are
     made in later rounds. \\
G6 & Adaptive operator selection (AOS) not examined in dynamic exchange.
   & Stateless UCB1 and contextual LinUCB select the mutation rate inside NSGA-III.
   & A negative result: under ILP guidance the two bandits gave the same outcomes and
     did not beat the best fixed rate. \\
\bottomrule
\end{tabularx}
\caption{Research gaps, our approach and current findings.}
\label{tab:gaps}
\end{table}

\section*{3. Model and methods}
\textbf{Dynamic pool.} In round $t$ the pool $V_t$ evolves as
$V_{t+1} = (V_t \setminus M_t \setminus X_t \setminus W_t) \cup A_t$: matched pairs
$M_t$ leave, patients exit to deceased-donor transplantation with probability $\mu_d$
($X_t$), patients who have waited $\tau$ months withdraw ($W_t$), and
$|A_t| \sim \mathrm{Poisson}(\lambda)$ new pairs arrive. Compatibility combines blood
type, a virtual crossmatch driven by cPRA, and transport time between six regional
nodes against a CIT limit. Cycles have at most 3 pairs and chains at most 4 transplants.
Every method faces the same patients through common random numbers.

\textbf{Objectives.} For a matching $M$ with patients $P$, highly sensitised patients
$S$ (cPRA $\geq 0.80$), and patients at distant ($R$) and central ($U$) nodes:
$F_1 = |M \cap P|$ (maximised); $F_2$ = Gini coefficient of waiting times after the
round; $F_3 = \bigl|\,|M \cap S|/|S| - |M \cap P|/|P|\,\bigr|$;
$F_4 = \bigl|\,|M \cap R|/|R| - |M \cap U|/|U|\,\bigr|$. $F_2$--$F_4$ are minimised.
The price of fairness in a round is
$\mathrm{PoF}_t = 100\,(F_1^{\mathrm{ILP}}(t) - F_1(t))/F_1^{\mathrm{ILP}}(t)$.

\textbf{Framework: ILP-guided NSGA-III.} An integer linear program (ILP) gives the
maximum number of transplants in each round. Its solution seeds NSGA-III, receives
priority in the repair step, and is kept in an external archive, so the Pareto front
always contains a volume-optimal matching. Rounds with at most 200 feasible matchings
are solved exactly by enumeration. A bandit (UCB1 or LinUCB) chooses the mutation rate
from the hypervolume gain. The implemented matching keeps at least 90\% of the maximum
volume (100\% and 80\% floors are also tested) and is closest to the ideal point in
$(F_2, F_3, F_4)$.

\textbf{Exact baseline: two-stage ILP.} Stage 1 finds the maximum volume $F_1^*$.
Stage 2 fixes $F_1 = F_1^*$ and minimises $F_2/r_2 + F_3/r_3 + F_4/r_4$, where each
$r_k$ is the objective's value in the volume-maximising matching. It returns the fairest
maximum matching and serves as a stronger baseline for the framework.

\textbf{Design.} Ten configurations, three withdrawal thresholds, 100 paired runs each,
with paired Wilcoxon tests and Bonferroni correction. The framework is compared with the
two-stage ILP at three levels: identical pools (one round), fixed trajectories (the same
pool in every round) and full trajectories (each method shapes its own future pools).
Sensitivity, scalability and out-of-distribution tests (PrefLib dataset 00036) complete
the evaluation.

\textbf{Modelling assumptions.} The exit rate, withdrawal thresholds, transport times,
the 24\,h CIT limit and the HLA proxy are modelling assumptions, not validated ANZKX
parameters. The arrival rate $\lambda = 10$ gives 65--70 simulated transplants per year,
somewhat above the reported 50--66, so we also test $\lambda = 7$ and $\lambda = 13$.

\section*{4. Key findings (simulation, 100 paired runs unless stated)}
\begin{itemize}
  \item \textbf{Fairness costs little over time.} With monthly clearing, the framework
        transplanted as many patients as the volume-maximising ILP over 25 rounds, while
        reducing the high-cPRA access gap $F_3$ by 21--24\%, the logistics gap $F_4$ by
        28--34\% and waiting-time inequality $F_2$ by 3--4\%. The mean PoF was 1.6--1.9\%
        per round (8.4\% in the first round, which clears the initial backlog, then
        flat), and 94--100\% of the per-round shortfall was recovered in later rounds.
  \item \textbf{Batching did not help.} With a threshold of 90 patients, batched clearing
        never cleared at $\tau = 6$ (65.6 highly sensitised patients withdrew per run,
        against 25.5 under monthly clearing) and lost 24\% and 12\% of transplants at
        $\tau = 12$ and 18. Better-tuned batching narrowed the gap but never closed it.
  \item \textbf{The exact two-stage ILP is stronger at maximum volume.} On 30 identical
        pools it reached $F_3 = 0.022$ and $F_4 = 0.035$, against 0.055 and 0.114 for the
        framework's maximum-volume choice. Over full trajectories it made 1.0--1.3 more
        transplants per run than the framework with its 90\% floor.
  \item \textbf{The front is useful below maximum volume.} On trajectory rounds the
        front contained solutions close to the exact reference at each volume level
        (at $\tau = 12$ and the 80\% level: 2.167 vs.\ 2.146 on the stage-2 score). An
        80\% floor lowered $F_3$ by 17--22\% relative to the two-stage ILP for
        2.1--2.9\% fewer transplants. The floor rule did not always pick the fairest of
        these solutions, which points to better selection rules.
  \item \textbf{ILP guidance, not the bandit, drives search quality.} Without guidance,
        NSGA-III never reached the volume optimum (0 of 30 pools; 32.6 vs.\ 42.3
        transplants) and its mean PoF was 28.3\%; with guidance it reached the optimum in
        30 of 30 pools and the PoF fell to 8.2\%.
  \item \textbf{Cost.} The two-stage ILP needed 5.08\,s per trajectory round against
        19.30\,s for the framework, but 233.3\,s on 75-vertex pools (67\% of solves
        time-limited) against 27.2\,s. The exact method returns one matching per volume
        level; the framework returns the whole front in one run.
  \item \textbf{PrefLib 00036} (12 real compatibility graphs with simulated arrivals,
        120 paired runs). Volume matched the ILP (75.03 vs.\ 75.38 transplants), the
        fairness gain was small ($F_3$ 0.060 vs.\ 0.062, not significant), and the
        two-stage ILP was again stronger at maximum volume.
\end{itemize}
The framework does not outperform every baseline. The two-stage ILP is the better tool
for the fairest matching at maximum volume, and a lexicographic rule that prioritises
highly sensitised patients gives them more transplants, at the cost of a larger $F_3$
gap. The framework's value lies in the trade-off solutions it offers when some volume may
be given up for fairness.

\section*{5. Open questions and opportunities for collaboration}
\begin{tabularx}{\textwidth}{@{}p{4.3cm}L@{}}
\toprule
Partner expertise & Joint work \\
\midrule
Market design, kidney exchange policy & Clearing frequency and batching rules that use
departure information (G4); comparison with policies in national programmes; dynamic
matching theory to explain why monthly clearing performed better here. \\
Fair division, computational social choice & Axiomatic analysis of the three fairness
criteria and of the selection rule; group versus individual waiting-time fairness;
fairness defined over time rather than per round (G5). \\
Multi-objective combinatorial optimisation & Exact or approximate Pareto fronts (for
example $\varepsilon$-constraint ILPs) to benchmark the evolutionary front below maximum
volume; selection rules aligned with the stage-2 objective; scaling beyond pools of a few
hundred vertices. \\
Healthcare operations, health economics & The deceased-donor pathway and logistics in
more realistic form (G2, G3); the value of a transplant and of waiting time; design of a
validation study with programme data, subject to ethics approval. \\
\bottomrule
\end{tabularx}

\section*{6. What we bring and possible next steps}
\textbf{What we bring.} A dynamic simulation of a kidney exchange pool with operational
constraints, informed by public ANZKX information; exact ILP and two-stage ILP
baselines; a tested multi-objective optimisation pipeline with paired-run statistics; and
open questions that sit between optimisation, market design and fairness.

\textbf{Possible joint work.} Depending on the partner's interests: regular online
meetings, a short research visit, a joint workshop and a joint paper. If there is mutual
interest, we would explore a small seed-funding scheme (for example the PSL Global Seed
Fund, with Adelaide University as the international partner) to start the collaboration,
and later a larger joint application.

\textbf{Contact:} Dr Yousef Amer, Adelaide University,
\href{mailto:yousef.amer@adelaide.edu.au}{yousef.amer@adelaide.edu.au}.

\end{document}
